\Needspace{12\baselineskip}
\originalchapter{流动失稳的机制}
\originalsection{界面扰动和Kelvin--Helmholtz失稳}
平稳流能否存在，与小扰动能否增长是两件事。将所有场写成基本态加小量，只保留小量的一阶项，再以 $\e^{\ii kx+\sigma t}$ 分解；$\Re\sigma>0$ 表示线性不稳定。这不会确定非线性饱和振幅，也不证明稳定基本态对任何有限扰动都稳定。

两层无黏、不可压流体密度 $\rho_1,\rho_2$，水平速度 $U_1,U_2$，界面张力 $\gamma_s$。各层无旋扰动势按 $\e^{-k|z|}$ 衰减，运动学边界把垂向速度与界面位移联系，动态边界给压力跳。消去振幅后，以波频 $\omega$ 表示的色散式为
\[
\left(\omega-k\frac{\rho_1U_1+\rho_2U_2}{\rho_1+\rho_2}\right)^2
=\frac{(\rho_2-\rho_1)gk+\gamma_s k^3}{\rho_1+\rho_2}
-\frac{\rho_1\rho_2k^2(U_1-U_2)^2}{(\rho_1+\rho_2)^2}.
\]
取下层密度 $\rho_2$，稳定分层重力项为正，剪切项为负。右端变负时频率出现虚部，流动剪切供能形成Kelvin--Helmholtz增长。无表面张力的跃变剪切模型在短波处可无界增长，反映连续剪切厚度被删去了，不能当成真实流体任意细尺度的结论。

\originalsection{连续剪切和Rayleigh判据}
基本速度 $U(y)$ 的二维无黏扰动流函数 $\psi=\phi(y)\e^{\ii k(x-ct)}$ 给Rayleigh方程 $(U-c)(\phi''-k^2\phi)-U''\phi=0$。乘以 $\phi^*/(U-c)$ 后积分，并取虚部，若 $\Im c>0$，则必须有 $\int U''|\phi|^2/|U-c|^2\dd y=0$。因此在相应刚性边界或衰减边界条件下，$U''$ 必须变号：拐点是无黏平行剪切失稳的必要条件。

这是必要条件而非充分条件。粘性Orr--Sommerfeld方程引入四阶导数和更多边界条件，某些无拐点流也会出现黏性不稳定。把Reynolds数增大后的临界值只用一条拐点判据取代，会漏掉实际机理。

\originalsection{扰动能量的用途}
将线性扰动动量与扰动速度点乘并积分，剪切生产项具有 $-\rho\int u'v'U'\dd V$ 的形式，黏性项为负耗散。扰动如何从基本流获得能量，需要相位使Reynolds应力与平均剪切配合。仅看速度扰动幅度大，不足以判断能量交换方向。

边界、重力、层化和黏性都可改变谱。源课的连续分层及黏性完整特征值问题需要专门分析，本卷给出方程和必要机制，不宣称完成所有波数与Reynolds数的稳定图计算。

\originalsection{界面色散式的逐步消元}
上层 $z>0$ 的扰动势 $\phi_1=A_1\e^{-kz}\e^{i(kx-\omega t)}$，下层 $\phi_2=A_2\e^{kz}\e^{i(kx-\omega t)}$，满足Laplace和远场衰减。界面 $\eta=\eta_0\e^{i(kx-\omega t)}$ 的运动学条件给 $-i(\omega-kU_i)\eta_0=\partial_z\phi_i|_0$，从而两势振幅具有相反法向符号。

线性Bernoulli压为 $p_i'=-\rho_i[(\partial_t+U_i\partial_x)\phi_i+g\eta]$。法向压跳含曲率，消去振幅得
\[
\rho_1(\omega-kU_1)^2+\rho_2(\omega-kU_2)^2
=(\rho_2-\rho_1)gk+\gamma_sk^3.
\]
配平方便是正文加权对流频率式。无剪切时稳定重力/毛细波；重流体在上使重力项负，Rayleigh--Taylor增长；同密度且无张力时任意剪切有KH增长。每个极限检验不同驱动的符号。

\originalsection{连续剪切、分层与Richardson条件}
二维无黏不可压扰动用流函数；把线性涡量方程写成 $(\partial_t+U\partial_x)\nabla^2\psi-U''\psi_x=0$，代入模式即得Rayleigh式。非平凡解的边界条件一般为墙面 $\phi=0$ 或无穷远衰减。虚部积分判据需边界项消失，不能略去这些前提。

Boussinesq稳定分层加入浮力方程后，得到Taylor--Goldstein式
\[
\phi''+\left[-k^2-\frac{U''}{U-c}+\frac{N_b^2}{(U-c)^2}\right]\phi=0.
\]
梯度Richardson数 $Ri=N_b^2/(U')^2$ 衡量分层恢复与剪切；在理想无黏、平行稳定分层及相应边界条件下，处处 $Ri\ge1/4$ 给经典充分稳定判据。$Ri<1/4$ 只取消这一保证，不证明必不稳定；分层、黏性和有限振幅问题要分别分析。

\originalsection{黏性Orr--Sommerfeld问题}
非维数以长度 $L$、速度 $U_0$，相同平行流扰动满足
\[
(D_y^2-k^2)^2\phi=ikRe[(U-c)(D_y^2-k^2)\phi-U''\phi].
\]
墙面无穿透和无滑移给 $\phi=\phi'=0$，四阶方程需要四个机械边界条件。$c$ 的虚部决定时间增长 $k\Im c$。$Re\to\infty$ 的阶数下降是奇异极限，边界层及 $U=c$ 附近临界层可能仍保留黏性，不能简单把四阶项归零后宣布全部谱一致。

在源课的目的中，这解释某些无拐点流为何仍可产生黏性不稳定。数值求谱需要处理壁边界、离散收敛与伪特征值，本卷没有实际生成全Re临界曲线，也不把写出特征方程当作已跑完谱图。

\originalsection{能量稳定与线性谱的不同结论}
扰动能量 $E'=\rho\int|\mathbf u'|^2/2$ 的变化为 $\dot E'=-\rho\int u'v'U'\dd V-2\mu\int e'_{ij}e'_{ij}\dd V$ 加适用的浮力项。利用不等式对生产项上界约束，可以证明某些参数下所有扰动能量衰减，这是比“所有本征值稳定”更强的结论。

非正规线性算子即使各本征值实部负，不同模叠加仍可能暂时放大；有限幅扰动也可能进入另一个吸引域。因此线性临界、能量稳定界与真实过渡Re不是同一数字。把三者分开能读懂源课程与后面非线性动力系统的联系。
